%% The Dungeon Master's Guide — pandoc/XeLaTeX template for the printable book (FR-004).
%%
%% This file is the reason the toolchain decision (01KXP7BCG6XFDHDRGZ57NXG90Q) was made.
%% pandoc + XeLaTeX was chosen over an HTML-to-PDF renderer, which would have reused the
%% site's stylesheet for free, BECAUSE the stakeholder reads this output on paper. What
%% is bought with that cost is below and nowhere else: real hyphenation, real
%% justification, widow and orphan control, a table of contents carrying page numbers,
%% and cross-references that resolve to a page.
%%
%% THE COST, stated plainly (plan.md § Complexity Tracking): this is a SECOND, INDEPENDENT
%% style system. src/theme/site/ (WP03) is the first. They can drift and nothing stops
%% them. Exact parity is not the requirement — "not embarrassing" is (NFR-004). The
%% shared decisions are declared once, immediately below, and referenced from both.
%%
%% Rendered by src/build/book.py. Do not run pandoc against this by hand and expect the
%% cross-references to resolve: book.py flattens them first (T018), and without that the
%% \pageref targets this file relies on are never written.

\documentclass[11pt,a4paper,twoside,openright]{book}

%% ---------------------------------------------------------------------------------
%% Shared decisions. THE ONE PLACE. If the site's palette or type changes, it changes
%% here too — this block and src/theme/site/'s stylesheet are the two halves of one
%% decision, and the only defence against the drift the plan accepted.
%% ---------------------------------------------------------------------------------
\usepackage{xcolor}
\definecolor{GuideInk}{HTML}{1A1A1A}      % body text — near-black, not black: kinder in print
\definecolor{GuideAccent}{HTML}{4051B5}   % Material's indigo, the site's primary
\definecolor{GuideMuted}{HTML}{6E6E6E}    % running heads, folios, rules
\definecolor{GuideRule}{HTML}{D8D8D8}

%% ---------------------------------------------------------------------------------
%% Fonts (T020). PINNED — see src/theme/book/fonts/README.md for the licences and the
%% reasoning. These names resolve inside the image book.py pins by digest
%% (pandoc/extra:3.8-ubuntu), which is the only place this template is supported. They
%% are NOT this laptop's fonts and NOT CI's system fonts; that is the point.
%%
%% DejaVu Sans Mono is not a default: it is the only mono in the image that carries the
%% box-drawing characters (U+2500, U+251C, U+2502, U+2514) the directory trees in
%% technical-suggestions.md and start.md are drawn with. Those live inside code fences,
%% where verbatim is verbatim and no substitution can reach them, so the font has to
%% carry them or the trees print as holes. Verified with fc-list inside the image;
%% src/build/book.py's COVERED_CHARACTERS records the result and enforces it.
%% ---------------------------------------------------------------------------------
\usepackage{etoolbox}
\usepackage{fontspec}
\defaultfontfeatures{Ligatures=TeX}
\setmainfont{DejaVu Serif}[
  Numbers        = OldStyle,
  Scale          = 0.92,        % DejaVu runs large; 11pt of it reads like 12
  SmallCapsFont  = {DejaVu Serif}
]
\setsansfont{DejaVu Sans}[Scale=0.92]
\setmonofont{DejaVu Sans Mono}[Scale=0.82]  % mono needs to shrink further to sit level

%% Emoji: NOT a font. src/build/book.py substitutes ✅/❌ for these dingbats before
%% XeLaTeX ever sees them — a deliberate decision (T019), reasoned in PRINT_SUBSTITUTIONS.
%% The book is printed; a colour emoji is grey ink anyway, and a dingbat sits on the
%% baseline with the text instead of riding above it.
\usepackage{pifont}

%% ---------------------------------------------------------------------------------
%% The page. A4, generous inner margin for the binding, 11pt body (FR-004).
%% Measure lands near 68 characters — the readable range, and the only reason the
%% margins are these numbers and not rounder ones.
%% ---------------------------------------------------------------------------------
\usepackage[
  a4paper,
  inner=28mm, outer=22mm,
  top=25mm,  bottom=28mm,
  headsep=8mm, footskip=14mm
]{geometry}

%% ---------------------------------------------------------------------------------
%% pandoc's own machinery. CALL THE PARTIAL — do not reimplement it.
%%
%% This one line is what supplies \tightlist (every list pandoc emits calls it),
%% \real and \newcounter{none} (uncaptioned longtables — i.e. every table in this
%% guide), \passthrough, \pandocbounded, the highlighting macros for fenced blocks, and
%% the paragraph-spacing setup. Hand-rolling these is how you find out at page 40 that
%% `\newcounter{none}` was the missing one, from an error pointing at generated code.
%%
%% It is placed HERE, before this template's own typography, so that everything below
%% wins: the partial sets \emergencystretch and parskip, and our decisions come after.
%%
%% Babel is deliberately not left to it. The partial keys off `lang`, and gets the
%% language from a \documentclass option this template does not have, so it would load
%% babel with no language at all. book.py passes no `lang`; British hyphenation is
%% requested explicitly below, and pdflang is set on hyperref by hand.
%% ---------------------------------------------------------------------------------
$common.latex()$

\usepackage{microtype}          % the margin kerning and glyph scaling LaTeX was chosen for
\usepackage{setspace}
\setstretch{1.10}

%% ---------------------------------------------------------------------------------
%% Widow and orphan control, and hyphenation (FR-004, T019).
%%
%% 10000 is "never", not "prefer not to": a single line of a paragraph stranded alone at
%% the top or foot of a page is the defect a reader notices without knowing why. LaTeX
%% pays for this by stretching the page, which \raggedbottom then lets it do honestly
%% instead of by opening up the leading.
%% ---------------------------------------------------------------------------------
\widowpenalty=10000
\clubpenalty=10000
\displaywidowpenalty=10000
\brokenpenalty=10000            % no hyphen dangling across a page break
\raggedbottom

\usepackage[british]{babel}     % British hyphenation patterns; the guide is written in
                                % British English ("satirise", "recognise"). Wrong
                                % patterns do not error — they just hyphenate wrongly,
                                % on every page, quietly.
\hyphenpenalty=200              % hyphenate rather than open rivers in a 68-char measure
\tolerance=1200
\emergencystretch=2em           % last resort before an overfull box; overrides the 3em
                                % the pandoc partial sets, which is loose for this measure

%% ---------------------------------------------------------------------------------
%% Tables. The guide's are wide — storytelling.md's good-tell/bad-tell comparison is two
%% columns of full sentences — and pandoc emits them as longtable, so they break across
%% pages by construction. \small buys the measure back.
%% ---------------------------------------------------------------------------------
%% longtable, booktabs, array, calc and \newcounter{none} all arrive with the partial.
\setlength{\LTpre}{\medskipamount}
\setlength{\LTpost}{\medskipamount}
\AtBeginEnvironment{longtable}{\small\setstretch{1.02}}
\renewcommand{\arraystretch}{1.25}

%% ---------------------------------------------------------------------------------
%% Block quotes. The guide leans on them heavily — the quick-read box in creator-kit.md
%% is one — so they get a rule rather than just an indent.
%% ---------------------------------------------------------------------------------
\usepackage{mdframed}
\renewenvironment{quote}{%
  \begin{mdframed}[
    leftline=true, rightline=false, topline=false, bottomline=false,
    linewidth=2pt, linecolor=GuideRule,
    innerleftmargin=10pt, innerrightmargin=0pt,
    innertopmargin=4pt, innerbottommargin=4pt,
    skipabove=\medskipamount, skipbelow=\medskipamount,
    backgroundcolor=white
  ]\small\itshape
}{%
  \end{mdframed}
}

%% ---------------------------------------------------------------------------------
%% Code. Fenced blocks are set in a tinted box; pandoc's highlighter fills in the rest.
%%
%% The highlighting-macros interpolation below is REQUIRED, not decoration: without it
%% every fenced block that names a language dies on an undefined \FunctionTok.
%%
%% WATCH OUT — a `%%` comment is a TeX comment, not a pandoc-template one. The template
%% engine expands variables inside comments too, so writing the variable's name in prose
%% here injects 30 lines of \newcommand into the middle of a sentence and the build dies
%% at "Missing \begin{document}" pointing at a line you never wrote. Learned the hard way.
%% ---------------------------------------------------------------------------------
\usepackage{framed}
\definecolor{shadecolor}{HTML}{F5F5F5}

%% \renew, not \new: the pandoc partial above already defines an empty `Shaded` and a
%% plain `Highlighting`. Defining them fresh fails with "Command \Shaded already
%% defined", pointing at a line of generated code. This block only exists to put the
%% tint and the size back.
\renewenvironment{Shaded}{\begin{snugshade}\small}{\end{snugshade}}
\DefineVerbatimEnvironment{Highlighting}{Verbatim}{commandchars=\\\{\},fontsize=\small}

\usepackage{enumitem}
\setlist{itemsep=2pt, parsep=2pt, topsep=4pt}

%% ---------------------------------------------------------------------------------
%% Headings.
%% ---------------------------------------------------------------------------------
\usepackage{titlesec}
\titleformat{\chapter}[display]
  {\normalfont\sffamily\huge\bfseries\color{GuideInk}}
  {\color{GuideAccent}\normalsize\sffamily\bfseries\MakeUppercase{\chaptertitlename\ \thechapter}}
  {12pt}{\Huge\raggedright}
\titlespacing*{\chapter}{0pt}{20pt}{28pt}
\titleformat{\section}{\normalfont\sffamily\Large\bfseries\color{GuideInk}}{\thesection}{0.6em}{}
\titleformat{\subsection}{\normalfont\sffamily\large\bfseries\color{GuideInk}}{\thesubsection}{0.6em}{}
\titleformat{\subsubsection}{\normalfont\sffamily\normalsize\bfseries\color{GuideMuted}}{}{0em}{}

%% ---------------------------------------------------------------------------------
%% Running heads and folios (FR-004).
%%
%% \nouppercase because the chapter titles are long sentences with em-dashes in them and
%% shouting them is worse than not having them. book.py issues \chaptermark with the
%% SHORT nav title right after each chapter opens, so the head says "How We Tell It" and
%% not "How we tell it — an alternate reality game, told diegetically", which does not
%% fit and would be silently overprinted into the margin.
%% ---------------------------------------------------------------------------------
\usepackage{fancyhdr}
\pagestyle{fancy}
\fancyhf{}
\renewcommand{\headrulewidth}{0.4pt}
\renewcommand{\footrulewidth}{0pt}
\fancyhead[LE]{\footnotesize\sffamily\color{GuideMuted}\nouppercase{\leftmark}}
\fancyhead[RO]{\footnotesize\sffamily\color{GuideMuted}\nouppercase{\rightmark}}
\fancyfoot[LE,RO]{\footnotesize\sffamily\color{GuideMuted}\thepage}
\fancypagestyle{plain}{\fancyhf{}\renewcommand{\headrulewidth}{0pt}%
  \fancyfoot[LE,RO]{\footnotesize\sffamily\color{GuideMuted}\thepage}}
\renewcommand{\chaptermark}[1]{\markboth{\thechapter.\ #1}{}}
\renewcommand{\sectionmark}[1]{\markright{#1}}

%% ---------------------------------------------------------------------------------
%% Links and cross-references.
%%
%% The \pageref targets that T018's flattening writes ("Chapter 6, Ethics (page 88)")
%% resolve against the \label pandoc puts on each chapter from its {#ch-ethics}
%% identifier. This is the mechanism the whole print-cross-reference story rests on;
%% hyperref must load for it, and it must load LAST.
%% ---------------------------------------------------------------------------------
\usepackage[unicode,hidelinks,pdfusetitle]{hyperref}
\hypersetup{
  pdftitle={$title$},
  pdfsubject={$subtitle$},
  pdflang={en-GB},
  bookmarksnumbered=true,
  bookmarksopen=true,
  linktoc=all
}
\urlstyle{same}

$for(header-includes)$
$header-includes$
$endfor$

\begin{document}

%% ---------------------------------------------------------------------------------
%% The title page (FR-004). Deliberately says what this object IS: not a book about the
%% guide, but the guide itself, printed. A reader who found it on a table should be able
%% to tell what they are holding and get back to the live version.
%% ---------------------------------------------------------------------------------
\frontmatter
\thispagestyle{empty}
\begin{titlepage}
\centering
\vspace*{0.16\textheight}
{\sffamily\bfseries\Huge\color{GuideInk} $title$\par}
\vspace{10mm}
{\color{GuideAccent}\rule{0.42\textwidth}{1.2pt}\par}
\vspace{10mm}
{\sffamily\large\color{GuideMuted} $subtitle$\par}
\vfill
{\small\itshape\color{GuideMuted}
The complete guide, cover to cover, for reading on paper.\par
\vspace{2mm}
Generated from the kit's own markdown --- nothing here is written twice.\par}
\vspace{6mm}
{\footnotesize\sffamily\color{GuideMuted} $if(base-url)$$base-url$$endif$\par}
\end{titlepage}

%% The table of contents. WITH PAGE NUMBERS (FR-004) — the book class's default, stated
%% here only because it is a named requirement and someone will check it.
%%
%% The number widths are not cosmetic. `book`'s defaults allow 1.5em for a section
%% number and 2.3em for a subsection, sized for "9.9" and "9.9.9". creator-kit.md has
%% fourteen sections, so its subsections reach "1.14" — wider than 2.3em, and LaTeX's
%% response is not to complain but to run the number straight into the title:
%% "1.109. Build your node". Widening is the fix; there is no warning to catch it by.
\makeatletter
\renewcommand*\l@section{\@dottedtocline{1}{0em}{2.6em}}
\renewcommand*\l@subsection{\@dottedtocline{2}{2.6em}{3.2em}}
\makeatother

\cleardoublepage
\pdfbookmark[0]{Contents}{toc}
\setcounter{tocdepth}{$if(toc-depth)$$toc-depth$$else$2$endif$}
\tableofcontents

\mainmatter

$body$

\end{document}
